\section{Conclusion and Future Directions}
\label{sec:conclusion}

In this paper, we have presented \textbf{SuiteMerge}, a systematic and
independent methodological audit of adaptive model-merging algorithms.
We have critically challenged
the prevailing evaluation protocols and exposed a severe, previously unreported
\textbf{Task Suite Bias} in the literature.
By partitioning a standard pool of four visual classification tasks into five
distinct multi-task evaluation suites, we demonstrated that the performance
rankings and claimed superiority of model-merging methods are highly sensitive
to domain distances and representational conflicts between tasks.
Under
realistic, correlated stream noise, unconstrained online Test-Time Adaptation
(AdaMerging) overfits catastrophically to local stream statistics, resulting in
representation collapse and severe accuracy degradation in high-conflict
regimes.
Conversely, we showed that Offline Few-Shot Validation Tuning
(\textbf{OFS-Tune}) using continuous linear trajectory constraints acts as a
robust analytical noise filter, consistently outperforming unconstrained online
TTA (AdaMerging) by up to \textbf{5.33\%} in simulation (and matching PolyMerge
within $1.02\%$), while successfully outperforming online PolyMerge by up to
\textbf{3.70\%} and online AdaMerging by up to \textbf{4.20\%} in real physical
weight-space deep networks, all while completely eliminating test-time
computational costs, latency, and privileged test-time routing assumptions.
Our work serves as a vital reality check for the model-merging community.
It
illustrates that evaluating algorithms on a single, fixed dataset suite can
obscure catastrophic local failures and lead to false progress.
For future work,
we urge researchers to shift away from monolithic benchmarks and adopt
multi-axial, multi-suite evaluation standards that explicitly analyze
representational conflicts across diverse task combinations.
Additionally, we advocate for scaling up our physical weight-space validation to modern large-scale benchmarks and foundation models as our \textbf{highest-priority next empirical step}.
Specifically, our immediate next work will apply SuiteMerge to instruction-tuned large language models (such as LLaMA-3 and Mistral-7B) across diverse, conflicting task relationships (e.g., reasoning vs.
coding vs.
mathematical problem solving using GLUE, MMLU, and HumanEval), as well as multi-modal Vision-Language Models (such as LLaVA-1.5) on conflicting visual question-answering and captioning tasks.
In these massive parameter spaces, the scale of representational conflict and structural depth are highly magnified.
Conducting this empirical transition is our top priority, as it will demonstrate the real-world scalability and engineering advantages of our proposed offline trajectory-tuning framework on frontier-scale architectures.
However, scaling static offline optimization like Nelder-Mead to
billion-parameter LLMs introduces significant computational and memory
challenges.
Standard derivative-free solvers require evaluating the full LLM's
multi-task performance at every iteration step, incurring high GPU memory
overhead and extreme evaluation latency across long validation streams.
To make
this roadmap highly actionable for large-scale foundation models, we propose
three concrete, easily implementable strategies utilizing existing libraries such
as Hugging Face PEFT, TRL, or Mergekit:
\begin{enumerate}
    \item \textbf{Parameter-Efficient Validation via Representative Subsets:} 
    Rather than evaluating model performance on extensive evaluation benchmarks
    per optimization step, the offline loss can be computed as cross-entropy or
    perplexity on a tiny, highly representative calibration subset (e.g., 10--20
    carefully selected instruction-response pairs per task, using prompt
    compression or key-token selection).
This dramatically reduces evaluation
    latency from minutes to seconds per step.
    \item \textbf{First-Order Coordinate Gradient Descent (OFS-Adam):} 
    To bypass the scaling limitations of Nelder-Mead in higher dimensions, we
    propose computing gradients of the validation loss directly with respect to
    the low-dimensional trajectory parameters $\theta_{k, j}$.
By registering
    the coefficients as differentiable PyTorch parameters within the merging
    graph, one can leverage standard backpropagation via PyTorch:
    \begin{equation}
    \label{eq:grad_ofs}
    \frac{\partial \mathcal{L}_{\text{val}}}{\partial \theta_{k, j}} = 
    \sum_{l=1}^L \frac{\partial \mathcal{L}_{\text{val}}}{\partial \alpha_k(l)} 
    \cdot \left( \frac{l}{L} \right)^j
    \end{equation}
    This formulation requires only standard few-shot first-order gradients and
    can be optimized within a few steps of Adam (OFS-Adam) without exhaustive
    local search, which is natively compatible with Hugging Face autograd
    pipelines.
    \item \textbf{Expert Parameter Offloading:} 
    To resolve the VRAM constraints of holding multiple massive LLM experts
    in GPU memory simultaneously during evaluation, we can leverage
    \texttt{accelerate}'s CPU offloading or device-map capabilities.
Expert
    parameters can be kept in host CPU memory and streamed sequentially into GPU
    memory layer-by-layer.
At each layer, the merged weight-space layer is
    computed on-the-fly, the forward activations are computed and cached, and the
    expert weights are immediately evicted from VRAM.
This maintains a constant GPU
    memory footprint equivalent to evaluating a single expert, making OFS-Tune
    highly practical on standard hardware.
\end{enumerate}
By adopting these engineering strategies, the model-merging community can scale
the robust, noise-filtering benefits of OFS-Tune to frontier foundation models,
achieving highly accurate, multi-task merged models without test-time
overfitting or deployment-time compute overhead.
Furthermore, we intend to empirically explore and build upon the mathematical formulations of alternative trajectory constraints introduced in Section~\ref{subsec:fidelity_mapping}---such as piecewise linear splines with attention-MLP boundaries and block-wise layer grouping---in future empirical studies.
This represents a promising path to determine if these formulations can successfully bridge the gap between robust noise-filtering and localized sensitivity spikes in modern Transformer architectures.
